%% kapa: exam=1 round=19 year=2008 period=1 form=A subject=accounting subject_ko=회계학 seq=38 number=118
%% source: Q-Net 감정평가사 기출문제 (cst00309) artlSeq=1001962 file=extracted/1008/1001962/2008년 제19회 감정평가사 1차시험 기출문제/1교시.hwp
%% conversion: kordoc 4.24.51 -> tools/gen1.py (auto); review=proofread-sonnet; flags=none
\begin{question}{118}
(주)한국은 종합원가계산시스템을 채택하고 있으며, 최초 공정은 A공정이며, 최종 공정은 B공정이다. 평균법을 적용한 20X1년 4월의 완성품환산량 단위당 원가는 다음과 같다.

{\small\centering
\begin{tabular}{lrlr}
\hline
직접재료원가 & ￦100/개 & 가공원가 & ￦200/개 \\
전공정(A공정)원가 & ￦500/개 & & \\
\hline
\end{tabular}
\par}

B공정에서 재료는 공정 50\% 시점에서 일괄 투입되며, 가공원가는 B공정 전체에 걸쳐 균등하게 발생한다. 4월말 B공정의 재공품은 1,000개이며, 가공원가의 완성도는 40\%이다. 이때 B공정의 기말재공품원가는? \par
\begin{choices}
  \item ￦480,000
  \item ￦580,000
  \item ￦680,000
  \item ￦780,000
  \item ￦880,000
\end{choices}
\end{question}
\answer{2}
